\section{Introduction}

% Paragraph 1: Introduce DeFi networks, credit exposures, risks, and measurement gaps

Decentralised finance (DeFi) has emerged as a sizeable and fast-evolving segment of the crypto-asset ecosystem, offering lending, trading, derivatives and payment-like services via smart contracts deployed across multiple blockchains \cite{Schar2021DeFi,Gogol2024SoKDeFi}. A defining feature of DeFi is composability: protocols routinely hold, lock, or accept tokens issued by other protocols as collateral, liquidity, or reserves. As a result, credit exposure is often implicit and token-mediated rather than expressed as bilateral contracts, creating a dense web of inter-protocol dependencies \cite{wu2025,Bertomeu2024}. When a token that is widely used as collateral or liquidity experiences a shock---through a price collapse, governance failure, exploit, or liquidation event---losses can propagate to protocols that rely on that token, potentially generating amplification and contagion across chains \cite{Aufiero2025,Gogol2024SoKDeFi}. Despite the policy relevance of these network effects, empirical work has been constrained by limited coverage and static snapshots, leaving a gap for dynamic, ecosystem-wide measurement that can track how exposure concentrations and transmission channels evolve over time \cite{wu2025}.

Policy bodies have repeatedly highlighted that DeFi can replicate traditional financial functions while inheriting familiar vulnerabilities, leverage, liquidity and maturity transformation, operational fragilities, and interconnectedness, often in forms that can unwind abruptly because key risk-management processes are automated in code \cite{FSB2023,IOSCO2023DeFi}. In particular, collateralised lending protocols rely on oracles and liquidation mechanisms that can trigger rapid deleveraging and fire-sale dynamics when prices move sharply or data feeds are disrupted \cite{DuleyEtAl2023Oracle,SasiBrodeskyNassr2023DeFiLiquidations,TianZhu2025Liquidations}. Stablecoins are a central node in these dynamics: they serve as settlement assets and collateral across protocols, and their de-pegs and runs can transmit stress through both on-chain and off-chain channels \cite{KosseEtAl2023Stablecoin,AnaduEtAl2023StablecoinsMMF,ESRB2025}. Yet regulators and researchers emphasise persistent data gaps that hinder monitoring and the construction of standardised risk indicators for the DeFi ecosystem \cite{FSB2023,IOSCO2023DeFi}. This measurement gap is increasingly salient as links between crypto markets and the broader financial system deepen \cite{ESRB2025}.

To address this gap, we train \textbf{DeXposure-FM}, a time-series, graph, and tabular foundation model for measuring and forecasting inter-protocol credit exposures in global decentralized financial networks. The model is trained on the \textbf{DeXposure} dataset, which reconstructs value-linked credit exposures and token flows across more than 4,300 protocols on 602 blockchains, covering over 24,300 tokens from 2020--2025 \cite{wu2025}. DeXposure operationalizes inter-protocol exposures by linking protocol balance-sheet dynamics (e.g., TVL) to token- and sector-level structure, enabling a unified view of volumes, topology, and compositional risk. DeXposure-FM follows the foundation-model paradigm: it pretrains a general representation on heterogeneous data and reuses it across downstream tasks, a strategy that has proven effective in time-series and structured-data settings \cite{liang2024fmts,eremeev2025g2tfm}. We train the model with Sharpness-Aware Minimization (SAM) to improve out-of-sample robustness by favoring parameter regions with uniformly low loss \cite{oikonomou2025sam}.

We evaluate DeXposure-FM on a suite of machine-learning benchmarks tailored to economic measurement and risk analytics. First, we perform multi-step forecasting of edge-level exposures and network-level statistics (e.g., concentration, density, and sector connectivity). Second, we conduct policy-relevant scenario analysis by generating shock-consistent paths of losses and contagion around major stress events (e.g., stablecoin de-pegs and exchange failures). Third, we study imputation and denoising of missing or noisy exposure and TVL histories using TVL-weighted metrics and downstream performance on forecasting and link-prediction tasks. Across tasks, we benchmark against classical econometric and latent-structure baselines (e.g., VAR and factor models) as well as temporal graph neural networks. %, following the evaluation philosophy established for DeXposure-style data.

On the financial economics front, DeXposure-FM provides a generative measurement layer for macroprudential monitoring in digital financial infrastructure. We use the trained model to construct dynamic measures of protocol-level systemic importance, sector-to-sector spillover indices, and early-warning indicators based on shifts in network concentration and dependence on fragile collateral. These AI-generated statistics support DeFi-specific stress testing and counterfactual policy evaluation (e.g., caps on bridge exposures or leverage, and alternative collateral rules), while clarifying when model-based measures are robust and when they require cautious interpretation by regulators and central banks. %More broadly, our contribution aligns with current priorities in AI-driven economic measurement, including evaluation on forecasting, policy analysis, and imputation \cite{nber2025ai}. 
By combining novel on-chain data with a multimodal foundation model, DeXposure-FM advances the measurement toolkit needed to understand systemic risk in DeFi and its interaction with the wider financial system \cite{naifar2025mapping,aufiero2025mapping}.